\section*{Discussion}

Our systematic comparison of population structure across SNVs, INDELs, and SVs in the 1000 Genomes high-coverage panel reveals a dual pattern: topological concordance overlaid with quantitative discordance that scales with variant size. All three variant classes recover the same continental population clusters, confirming that the dominant axes of human genetic variation reflect shared demographic history detectable regardless of mutational mechanism. However, the magnitude of population differentiation is systematically attenuated for larger variants, and a convergence analysis establishes that the discordance between SNV-based and SV-based population structure is not a statistical artifact of low variant count but rather an irreducible biological signal.

\subsection*{A selection-mediated differentiation gradient}

The observed size-dependent $\fstmath$ gradient---from 0.075 for 1~bp INDELs through 0.060 for SVs---is consistent with differential purifying selection, though alternative explanations cannot be excluded without formal testing. Larger genomic rearrangements disrupt more sequence per event and are more likely to affect coding regions, regulatory elements, or chromatin organization, leading to stronger negative fitness effects \citep{abel2020, sudmant2015}. Under purifying selection, deleterious alleles are removed before they can drift to appreciable frequency differences between populations, compressing the $\fstmath$ distribution. The continuity of the gradient across the conventional INDEL-SV boundary at 50~bp is particularly informative: it suggests that the distinction between these variant classes in population-genetic behavior is quantitative, not qualitative. Rather than two discrete categories with distinct evolutionary dynamics, INDELs and SVs appear to represent a mutational continuum along which the strength of purifying selection increases smoothly with variant size.

The corroborating evidence from the site frequency spectrum supports this interpretation. SVs are dramatically enriched for rare variants (59--63\% with MAF 0.01--0.05, compared to 18--32\% for SNVs under uniform filtering), consistent with the expectation that strongly selected variants are maintained at low frequencies by mutation-selection balance. The SFS skew survives harmonization of MAF thresholds across variant types, arguing against ascertainment bias as the primary driver. However, we emphasize that mutation-rate heterogeneity across variant size classes could also produce a size-dependent $\fstmath$ gradient under neutrality: if larger variants have lower per-generation mutation rates, the resulting younger allele age distribution would mimic the effects of purifying selection. Distinguishing between selection and mutation-rate explanations requires forward simulations calibrated to empirical mutation rates for each size class, which we leave for future work.

\subsection*{Insertion-deletion asymmetry: a challenge to symmetric models}

The 25\% reduction in $\fstmath$ for insertions relative to deletions of comparable size (0.046 versus 0.062) is unexpected under models that treat insertions and deletions symmetrically. Several mechanisms could underlie this asymmetry. Insertions add novel sequence to the genome, which may be more likely to disrupt reading frames through frameshift, create cryptic splice sites, or interfere with regulatory grammar than deletions that simply remove sequence. Additionally, the reference genome itself may be biased: if the human reference preferentially represents ancestral states \citep{schneider2017}, then insertions relative to the reference may correspond disproportionately to derived alleles, which are on average younger and rarer due to the combined effects of drift and selection. Distinguishing between these explanations will require functional annotation of insertion and deletion variants, which was beyond the scope of this study but represents a natural extension. We also note a technical caveat: short-read SV calling is known to under-detect insertions relative to deletions \citep{ho2020}, and the 3:1 DEL-to-INS ratio (14,091 versus 4,428) likely reflects both biology and ascertainment bias. Insertions that pass short-read-based calling filters may represent a biased subset---potentially enriched for variants in less complex genomic regions that happen to be under weaker selective constraint---which could inflate the apparent strength of purifying selection on the detected insertion set. Long-read-based SV catalogs, which recover more balanced insertion-to-deletion ratios, will be essential to validate this finding.

\subsection*{The convergence plateau as a diagnostic tool}

The convergence analysis we developed provides a framework that may generalize for decomposing cross-type concordance into power and biology components, though its broader applicability remains to be tested beyond this study. The key insight is that the rate and ceiling of convergence differ qualitatively between variant pairs that share the same underlying signal (SNV-INDEL, ceiling $\sim$0.99) and those that do not (SNV-SV, ceiling $\sim$0.60). This framework could be applied broadly---for example, to compare SNP-array-based versus whole-genome-based population structure, to evaluate how different SV callers or sequencing technologies affect population-genetic inference, or to assess concordance between genetic and epigenetic markers of population differentiation. The count-matched comparison at 18,519 variants, which showed SNV-INDEL concordance at $r \approx 0.89$ versus SNV-SV at $r \approx 0.55$, provides an intermediate data point: at the SV count, SNV-INDEL concordance is already recovering, while SNV-SV concordance is near its final plateau. The convergence curve thus serves as a diagnostic: a plateau well below 1.0 signals genuine biological discordance, while a trajectory still rising at the available count signals that more markers would improve concordance.

\subsection*{Implications for population genomics with structural variants}

Our findings have several practical implications. First, the near-perfect $\fstmath$ rank-order correlation ($> 0.997$) across variant types confirms that \emph{which} populations are most differentiated is robust to variant class. Ancestry inference tools calibrated on SNVs should correctly rank populations regardless of the variant type used, validating current practice. Second, the 28\% systematic attenuation of SV $\fstmath$ relative to SNV $\fstmath$ means that selection scans based on $\fstmath$ outliers will have different power depending on variant type: a given $\fstmath$ percentile threshold will capture different genomic targets in SNVs versus SVs, and direct comparison of $\fstmath$ values across variant types requires calibration. Third, the insertion-deletion asymmetry suggests that SV burden analyses that treat insertions and deletions interchangeably may obscure meaningful differences in their population-genetic behavior.

\subsection*{Limitations and future directions}

Several limitations should be acknowledged. The SV call set from the 1kGP high-coverage panel is based on short-read sequencing and is therefore incomplete, particularly for complex SVs, large insertions, and variants in repetitive regions \citep{ho2020}. The size range of SVs in our analysis (predominantly 50--200~bp) is narrower than the full SV spectrum; long-read-based call sets may reveal different patterns for larger variants. We did not perform functional annotation of variants, which would strengthen the mechanistic interpretation by directly linking high-$\fstmath$ variants to specific functional categories. No formal demographic modeling was conducted to assess whether the observed $\fstmath$ gradient is expected under neutrality with mutation-rate heterogeneity, or whether purifying selection is required. ADMIXTURE analysis, which could reveal type-specific differences in ancestry proportions for admixed populations, was not performed. Bootstrap confidence intervals were not computed for the main statistics; however, the consistency of patterns across 325 population pairs, five super-populations, and three random seeds provides empirical evidence of robustness. The core $\fstmath$ and Procrustes analyses were conducted with the original type-specific MAF thresholds (0.01 for SNVs/INDELs, 0.005 for SVs); while the SFS analysis was repeated with uniform MAF $\geq 0.01$ filtering confirming robust SFS differences, ideally the $\fstmath$ gradient should also be verified under harmonized thresholds to fully exclude ascertainment-driven effects. Additionally, no sensitivity analysis was performed removing related individuals from the 602 trios in the 1kGP panel; while relatedness has minimal impact on allele frequency estimation in large samples, its effect on cross-type Procrustes concordance specifically has not been tested. Finally, our analyses are based on a single dataset; replication in independent cohorts (HGDP, gnomAD-SV) and with long-read-based SV call sets would strengthen generalizability.

Future work should prioritize functional enrichment analysis of population-differentiated variants across types, formal demographic inference per variant class, and extension to long-read SV catalogs that capture the full size spectrum of structural variation. The convergence analysis framework developed here provides a template for these comparisons.
